%% ============================================================================
%%  applications/template.tex  —  SKELETON RESUME (single source template)
%% ----------------------------------------------------------------------------
%%  This is the ONLY resume template. The orchestrator copies this file into a
%%  position folder, then the writer fills the Experience / Projects / Skills
%%  sections from the context.md bullet pool. It replaces the old per-track
%%  templates (full-stack / ai-ml / dev-ops) — track is now a doctrine-routing
%%  label, not a file choice, so any track (including robotics/autonomy) fills
%%  this same skeleton.
%%
%%  WHAT IS FIXED (do not change): header (\name, \contact) and Education.
%%  WHAT GETS FILLED: Experience entries, Projects entries, Technical Skills.
%%
%%  rSectionEntry signature — 5 positional params, ALWAYS supply all 5 braces:
%%     {1 name}{2 dates}{3 tagline/role/link}{4 technologies}{5 location}
%%     - param 4 (technologies) renders after the name with a " | " separator;
%%       leave it {} for experiences, fill it for projects (the tech-stack line).
%%     - any unused param is an empty {}.
%%  rSet signature — {1 bucket label}{2 comma-separated items}.
%%
%%  PATH NOTE: this template sits at resume/applications/template.tex, beside
%%  template.cls, so the class path is just `template`. On copy into
%%  applications/{year}/{company}/{role}/, the orchestrator rewrites it to
%%  ../../../template (three levels up to applications/, where template.cls lives).
%% ============================================================================

\documentclass[
	11pt
]{../../../template}

\name{Vedant Desai}
\contact{(248) $\hspace{-0.1em}$ 704 $\hspace{-0.4em}$ -- $\hspace{-0.4em}$ 4852 \\ vedantde@umich.edu \\ \href{https://github.com/Verdent06}{github.com/Verdent06} \\ \href{https://linkedin.com/in/vedantde06}{linkedin.com/in/vedantde06}}

\begin{document}

	% ---- EDUCATION (fixed; writer does not modify) --------------------------
	\begin{rSection}{E}{ducation}
		\begin{rSectionEntry}
			{University of Michigan}
			{Expected May 2028}
			{B.S. in Computer Science and Economics}
			{}
			{Ann Arbor, MI}

			\item \textbf{GPA:} 3.66 / 4.0
			\item \textbf{Coursework:} Data Structures \& Algorithms, Intro to Statistics and Data Analysis, Microeconomics, Macroeconomics, Discrete Mathematics, Calculus III, Physics (Mechanics)
		\end{rSectionEntry}
	\end{rSection}

	% ---- EXPERIENCE (writer fills from context.md pool) ---------------------
	% Order by impact density for THIS role, not chronology (see resume.md I.3).
	% First bullet of each entry is the hook: lead with the problem/outcome,
	% not the technology (see resume.md I.4). 2-3 bullets per experience.
	\begin{rSection}{E}{xperience}
		\begin{rSectionEntry}
			{Michigan Data Consulting (MDC)}
			{Jan 2026 -- May 2026}
			{Data Engineer --- Michigan Campaign Finance Network}
			{}
			{Ann Arbor, MI}

			\item Replaced manual Michigan campaign-finance research --- portal searches capped by the Bureau of Elections, irregular Excel exports, hand normalization at $\sim$2 hours per committee --- with a Requests + Pandas ETL that ingests filings directly, eliminating $\sim$800 hours of manual pulls across 400 tracked PACs.
			\item Delivered a production Flask REST API on AWS EC2 as the sole engineer on a 5-month MCFN contract, wiring ingested data and PAC rankings into the nonprofit's public-facing research workflow.
		\end{rSectionEntry}

		\begin{rSectionEntry}
			{Vylet}
			{May 2026 -- Present}
			{Founder --- Automated lead-sourcing platform for PE/search-fund; live product generating \$1,500 MRR across three paying clients}
			{Python, LangGraph, Redis, Celery, Docker}
			{\href{https://vyletdata.com}{vyletdata.com}}

			\item Launched Vylet, automating a $\sim$30-minute manual process per business into a Dockerized LangGraph pipeline generating 30 scored leads in 30 minutes --- a 30x speedup --- with Redis/Celery workers on a recurring cycle.
			\item Diagnosed a name-collision defect in ownership-verification logic that was incorrectly rejecting valid acquisition targets sharing a name with an unrelated business elsewhere; the fix lifted the pipeline's lead-qualification rate from 79\% to 89\% with zero change to sourcing volume.
		\end{rSectionEntry}
	\end{rSection}

	% ---- PROJECTS (writer fills from context.md pool) -----------------------
	% 1-2 bullets each. Tech-stack line goes in param 4 (derived from the bullets
	% selected for this project). Optional live link goes in param 3 as \href.
	\begin{rSection}{P}{rojects}
		\begin{rSectionEntry}
			{Granular Synthesizer Plugin}
			{\href{https://github.com/Verdent06/granular-synth}{github.com/Verdent06/granular-synth}}
			{Real-time audio DSP plugin built from scratch in C++/JUCE --- sine oscillator through full granular engine with lock-free threading and release-ready VST3/AU binaries}
			{C++, JUCE}
			{}

			\item Enforced the audio thread's zero-allocation constraint by pre-allocating a MemoryPool\textless{}Grain, 64\textgreater{} slab per voice at plugin startup --- a fixed free-list that acquires and releases grain slots without touching the heap --- so processBlock() never calls new or delete after prepareToPlay() completes.
			\item Wired UI-to-audio delivery through a hand-rolled SPSC FIFO (64-slot fixed array, atomic head/tail with acquire/release ordering) so slider changes reach the audio thread without a mutex --- and routed WAV handoff via atomic shared\_ptr swap so the UI thread never blocks processBlock().
			\item Configured CMake to compile VST3 and AU bundles from a single JUCE codebase --- targeting a macOS universal binary (arm64 + x86\_64) for native Apple Silicon and Intel --- and audited every processBlock() path against a real-time safety checklist confirming zero heap allocations and zero lock acquisitions before release builds.
		\end{rSectionEntry}

		\begin{rSectionEntry}
			{SignalWeaver}
			{}
			{Multi-signal financial research platform combining fundamentals, macro indicators, and news sentiment (research assistant; not investment advice)}
			{Python, PostgreSQL, pgvector, Docker, GitHub Actions}
			{}

			\item Implemented semantic search over stored financial news with pgvector cosine similarity (768-d MPNet embeddings), hitting 49ms p50 / 99ms p99 over 90 queries returning top-5 or top-10 matches.
			\item Built a React/TypeScript dashboard for composite scores and fundamental/sentiment/macro breakdowns; used it to analyze 90 distinct tickers with scores persisted to Postgres for history views.
			\item Containerized the stack with Docker Compose (API + Postgres/pgvector + nginx frontend) and a GitHub Actions CI pipeline (frontend build, pytest, API image build on main).
		\end{rSectionEntry}
	\end{rSection}

	% ---- TECHNICAL SKILLS (writer picks 3-5 buckets from context.md) --------
	% ~80 char limit per line after the label. Trim tail items to fit; promote a
	% tail item only if it is a hard ATS keyword from the JD (see resume.md I.5).
	\begin{rSection}{T}{echnical Skills}
		\begin{rSet}{Languages}{Python, C++, SQL}
		\end{rSet}
		\begin{rSet}{Frameworks}{Flask, JUCE}
		\end{rSet}
		\begin{rSet}{Databases}{PostgreSQL, Redis, pgvector}
		\end{rSet}
	\end{rSection}

\end{document}
